\section{Research Experience}

\entry{PhD Researcher, Byzantine-Resilient Distributed Learning}{Nov 2026--present}
  {CMAP, \textnormal{\'E}cole Polytechnique, France \textnormal{\textbar} ATLAS chair}
Research topic: trade-offs between adversarial robustness, data ownership
verification and efficient learning. Initial project on the Byzantine fault
tolerance of signSGD.\par
Supervisor: El-Mahdi El-Mhamdi. Co-directors: Katia Meziani and Vladimir Kostic.
Co-advisors: Peva Blanchard and S\'ebastien Rouault.

\entry{Research Intern, Byzantine-Resilient Distributed Learning}{Apr--Oct 2026}
  {CMAP, \textnormal{\'E}cole Polytechnique, France \textnormal{\textbar} ATLAS chair}
Supervised by El-Mahdi El-Mhamdi, with Peva Blanchard and S\'ebastien Rouault.
\begin{itemize}
  \item Co-developed \href{https://links.arthurdanjou.fr/krum}{\textbf{Krum}}, an
    open-source Python/PyTorch library for evaluating distributed learning under
    Byzantine attacks, released on PyPI under the MIT licence.
  \item Implemented gradient aggregation rules, a zero-copy model wrapper,
    neural network architectures and centralised training simulations.
  \item Reproduced key experimental behaviours from NeurIPS 2017, ICML 2018
    and ICML 2023 protocols, with documented adaptations, including a
    reduced-scale decentralised MoNNA case study.
  \item Wrote component tests, tutorials and reference documentation, and
    configured continuous integration and software release workflows.
  \item Compared related research libraries and contributed to the companion
    manuscript describing Krum and its experimental methodology.
\end{itemize}
